% Source printed pp. 122-129, physical PDF pp. 130-137.
\section{基本性质}\label{src:4.3-section}
\begin{epigraph}
``当然, 这里我研究的是模叠, 而不是模空间. 对于不熟悉叠的读者, 不必担心: 基本上, 这只是说我可以假装模空间是光滑的, 并且它上面存在万有族.''

谁没有在报告中听过这些话或类似的说法? 又有谁没有想象过揪住报告人的衣领, 摇晃他, 直到他说清楚这些话究竟是什么意思?
\epigraphauthor{Joe Harris 和 Ian Morrison \cite[第 $140$ 页]{HM98}}
\end{epigraph}
我们定义代数叠态射的一些基本性质 (例如闭浸入, 仿射, 局部有限型, 光滑, 拟分离, 拟有限和平展), 以及代数叠的一些基本性质 (例如局部 Noether, 拟紧, 连通, 约化, 正则和 Noether). 我们还引入代数叠 $\mathcal X$ 的底层拓扑空间 $|\mathcal X|$.
\subsection{态射的性质}\label{src:4.3.1-subsection}
回忆, $\Sch_{\text{ét}}$ 上预叠的态射 $\mathcal X\to\mathcal Y$ 可由概形表出 (相应地, 可表), 是指对于每个来自概形的态射 $T\to\mathcal Y$, 基变换 $\mathcal X\times_{\mathcal Y}T$ 是概形 (相应地, 代数空间); 见定义~\bookref{4.1.1}~和~\bookref{4.1.3}. 这两个概念显然在基变换下稳定. 可由概形表出的态射显然也在复合下稳定, 下述引理说明可表态射同样如此.
\begin{lemma}\label{src:4.3.1}
\begin{enumerate}[(1)]
\item 若 $\mathcal X\to\mathcal Y$ 是 $\Sch_{\text{ét}}$ 上预叠的可表态射, 且 $\mathcal Y$ 是代数空间, 则 $\mathcal X$ 是代数空间.
\item 可表态射的复合可表.
\end{enumerate}
\end{lemma}
\begin{proof}
对于第一条, 若 $V\to\mathcal Y$ 是来自概形 $V$ 的平展呈示, 则基变换 $\mathcal X_V$ 是代数空间. 由于对角线 $\mathcal X\to\mathcal X\times_{\mathcal Y}\mathcal X$ 沿 $V\to\mathcal Y$ 作基变换后成为单态射 $\mathcal X_V\to\mathcal X_V\times_V\mathcal X_V$, 平展下降蕴涵 $\mathcal X\to\mathcal X\times_{\mathcal Y}\mathcal X$ 也是单态射. 因而 $\mathcal X$ 等价于层 (习题~\bookref{3.4.39}). 此外, 基变换 $\mathcal X_V\to\mathcal X$ 是代数空间的态射, 平展, 满射且可由概形表出. 取平展呈示 $U\to\mathcal X_V$, 则复合 $U\to\mathcal X_V\to\mathcal X$ 平展, 满射且可由概形表出, 从而 $\mathcal X$ 是代数空间. 第二条由第一条得到.
\end{proof}
% Keep the definition's introduction together with the first defining condition.
\newpage
\begin{definition}\label{src:4.3.2}
设 $\mathcal P$ 是概形态射的性质.
\begin{enumerate}[(1)]
\item 若对于每个概形的平展满射 $X'\to X$, 态射 $X\to Y$ 满足 $\mathcal P$ 当且仅当 $X'\to X\to Y$ 满足 $\mathcal P$, 则称 $\mathcal P$ \newterm{在源上平展局部}{étale local on the source}. 若对于每个概形的平展满射 $Y'\to Y$, 态射 $X\to Y$ 满足 $\mathcal P$ 当且仅当 $X\times_Y Y'\to Y$ 满足 $\mathcal P$, 则称 $\mathcal P$ \newterm{在靶上平展局部}{étale local on the target}. 在源/靶上光滑局部 (相应地, fppf 局部, fpqc 局部) 的概念类似定义.
\item 若 $\mathcal P$ 在源和靶上平展局部 (相应地, 光滑局部), 且在复合和基变换下稳定, 则称 Deligne--Mumford 叠 (相应地, 代数叠) 的态射 $\mathcal X\to\mathcal Y$ 具有性质 $\mathcal P$, 如果对于所有给出图
\[
\begin{tikzcd}
U\arrow[r]&\mathcal X\times_{\mathcal Y}V\arrow[r]\arrow[d]\arrow[dr,phantom,"\square"]&V\arrow[d]\\
{}&\mathcal X\arrow[r]&\mathcal Y,
\end{tikzcd}
\]
的平展 (相应地, 光滑) 呈示 $V\to\mathcal Y$ 和 $U\to\mathcal X\times_{\mathcal Y}V$, 复合 $U\to V$ 都具有 $\mathcal P$. 只需对某些特定的呈示 $V\to\mathcal Y$ 和 $U\to\mathcal X\times_{\mathcal Y}V$ 检查此性质.
\item 可由概形表出的代数叠态射 $\mathcal X\to\mathcal Y$ 具有性质 $\mathcal P$, 是指对于每个来自概形的态射 $T\to\mathcal Y$, 基变换 $\mathcal X\times_{\mathcal Y}T\to T$ 具有 $\mathcal P$. 特别地, 代数叠的态射 $\mathcal X\to\mathcal Y$ 是同构, 开浸入, 闭浸入, 局部闭浸入, 仿射或拟仿射, 是指它可由概形表出并具有相应的性质.
\end{enumerate}
\end{definition}
平坦, 光滑 (相应地, 相对维数为 $n$ 的光滑), 满射, 局部有限表现和局部有限型等性质在源和靶上光滑局部. 利用定义~\bookref{4.3.2}(2), 这些性质推广到代数叠的态射. 同样, 平展性和非分歧性在源和靶上平展局部, 从而推广到 Deligne--Mumford 叠的态射. 这些性质在复合和基变换下稳定.
\begin{example}\label{src:4.3.3}
若 $G\to S$ 是作用于代数空间 $U\to S$ 的光滑仿射群概形, 则 $[U/G]\to S$ 平坦 (相应地, 光滑, 满射, 局部有限表现, 局部有限型) 当且仅当 $U\to S$ 如此. 特别地, 利用定理~\bookref{4.1.17}~和~\bookref{4.1.23}~证明中的商叠呈示, 可得 $\mathcal M_g$ 在 $\mathbb Z$ 上局部有限型, 且 $\Coh_{r,d}(C)$ 和 $\cBun_{r,d}(C)$ 都在 $k$ 上局部有限型.
\end{example}
\begin{exercise}[简单]\label{src:4.3.4}
令 $G=\mathbb Z/2$ 通过 $(-1)\cdot(x,y)=(-x,-y)$ 作用于域 $k$ 上的 $\mathbb A^2$. 证明 $[\mathbb A^2/G]$ 光滑, 而概形商 $\mathbb A^2/G:=\Spec k[x,y]^G$ 奇异.
\end{exercise}
\begin{exercise}[简单]\label{src:4.3.5}
证明代数叠态射 $\mathcal X\to\mathcal Y$ 的对角线局部有限型.
\end{exercise}
现在证明, 代数叠的某些态射在靶上光滑局部. 它们甚至在靶上 fppf 局部, 但这一点推迟到习题~\bookref{7.3.5}~讨论.
\begin{proposition}\label{src:4.3.6}
设 $\mathcal P$ 是代数叠态射的下列性质之一: 可表, 同构, 开浸入, 闭浸入, 仿射或拟仿射. 考虑代数叠的 Cartesian 图
\[
\begin{tikzcd}
\mathcal X'\arrow[r]\arrow[d]\arrow[dr,phantom,"\square"]&\mathcal Y'\arrow[d]\\
\mathcal X\arrow[r]&\mathcal Y,
\end{tikzcd}
\]
其中 $\mathcal Y'\to\mathcal Y$ 光滑且满射. 则 $\mathcal X\to\mathcal Y$ 具有 $\mathcal P$ 当且仅当 $\mathcal X'\to\mathcal Y'$ 具有 $\mathcal P$.
\end{proposition}
\begin{proof}
我们证明 $\Leftarrow$ 方向, 因为另一方向显然. 对于可表性, 可以假设 $\mathcal Y$ 和 $\mathcal Y'$ 是概形. 只需证明 $\mathcal X$ 中位于概形 $T$ 上的对象 $a$ 的每个自同构 $\alpha:a\to a$ 都平凡. 将 $a:T\to\mathcal X$ 沿 $\mathcal X'\to\mathcal X$ 作基变换得到的 $T'$ 等同于 $T\times_{\mathcal Y}\mathcal Y'$, 从而是 $T$ 上的光滑概形. 由于光滑态射在平展局部具有截面 (推论~\bookref{A.3.5}), 存在通过 $T'$ 分解的平展覆盖 $g:\widetilde T\to T$. 自同构 $\alpha$ 定义了 $\Aut_T(a)$ 在 $T$ 上的截面. 由于 $\Aut_T(a)$ 是 $(\Sch/T)_{\text{ét}}$ 上的层且 $g^*\alpha=\id$, 有 $\alpha=\id$.

对于其他性质, 已知 $\mathcal X\to\mathcal Y$ 可表, 因而只需假设 $\mathcal Y$, $\mathcal Y'$ 和 $\mathcal X'$ 是概形, $\mathcal X$ 是代数空间. 此时可以应用 fppf 层成为概形的下降判据 (命题~\bookref{3.3.17}), 得出 $\mathcal X$ 是概形. 于是结论成立, 因为概形态射的这些性质是在靶上 fpqc 局部的性质 (\bookref{3.1.26}).
\end{proof}
\subsection{代数空间与代数叠的性质}\label{src:4.3.2-subsection}
\begin{definition}[代数空间与代数叠的性质]\label{src:4.3.7}
设 $\mathcal P$ 是概形的平展局部 (相应地, 光滑局部) 性质, 即若 $X\to Y$ 是概形的平展 (相应地, 光滑) 满射, 则 $X$ 具有 $\mathcal P$ 当且仅当 $Y$ 具有 $\mathcal P$. 称 Deligne--Mumford 叠 (相应地, 代数叠) $\mathcal X$ 具有性质 $\mathcal P$, 如果对于某个平展 (相应地, 光滑) 呈示 $U\to\mathcal X$ (等价地, 对所有这样的呈示), 概形 $U$ 具有 $\mathcal P$.
\end{definition}
局部 Noether, 约化和正则等性质是光滑局部的 (命题~\bookref{3.1.21}).
\begin{example}\label{src:4.3.8}
设 $G\to S$ 是作用于 $S$ 上概形 $U$ 的光滑仿射群概形. 则 $[U/G]$ 局部 Noether, 约化或正则, 当且仅当 $U$ 如此.
\end{example}
\begin{definition}[子叠]\label{src:4.3.9}
若 $\mathcal X$ 是代数叠, 则子叠 $\mathcal Z\subseteq\mathcal X$ 是闭的 (相应地, 开的, 局部闭的), 是指诱导的态射 $\mathcal Z\to\mathcal X$ 是闭浸入 (相应地, 开浸入, 局部闭浸入).
\end{definition}
\begin{exercise}\label{src:4.3.10}
对于光滑仿射群概形 $G\to S$ 在 $S$ 上概形 $U$ 上的作用, 证明 $[U/G]$ 的闭 (相应地, 开) 子叠与 $U$ 的 $G$-不变闭 (相应地, 开) 子概形之间存在等价.
\end{exercise}
\subsection{分离性}\label{src:4.3.3-subsection}
代数叠的分离性通过对角线定义.
\begin{definition}\label{src:4.3.11}
\begin{enumerate}[(1)]
\item 代数叠态射 $\mathcal X\to\mathcal Y$ 具有\newterm{仿射对角线}{affine diagonal} (相应地, \newterm{拟仿射对角线}{quasi-affine diagonal}, \newterm{分离对角线}{separated diagonal}), 是指对角线 $\mathcal X\to\mathcal X\times_{\mathcal Y}\mathcal X$ 仿射 (相应地, 拟仿射, 分离). 代数叠 $\mathcal X$ 具有仿射对角线 (相应地, 拟仿射对角线, 分离对角线), 是指 $\mathcal X\to\Spec\mathbb Z$ 如此.
\item 代数叠态射 $\mathcal X\to\mathcal Y$ 拟分离, 是指对角线 $\mathcal X\to\mathcal X\times_{\mathcal Y}\mathcal X$ 和第二对角线 $\mathcal X\to\mathcal X\times_{\mathcal X\times_{\mathcal Y}\mathcal X}\mathcal X$ 都拟紧. 代数叠 $\mathcal X$ 拟分离, 是指它在 $\Spec\mathbb Z$ 上拟分离.
\item 代数叠态射 $\mathcal X\to\mathcal Y$ 有限表现, 是指它局部有限表现, 拟紧且拟分离.
\item 代数叠的可表态射 $\mathcal X\to\mathcal Y$ 分离, 是指态射 $\mathcal X\to\mathcal X\times_{\mathcal Y}\mathcal X$ 是固有的, 此态射可由概形表出 (习题~\bookref{4.2.4}). 一般情形的分离态射在定义~\bookref{4.8.1}~中定义.
\end{enumerate}
\end{definition}
对角线上的条件转化为 Isom 层上的条件, 因为 $\mathcal X\to\mathcal X\times\mathcal X$ 沿来自概形 $S$ 的态射 $(a,b):S\to\mathcal X\times\mathcal X$ 作基变换后等同于 $\underline{\Isom}_{\mathcal X(S)}(a,b)$ (见习题~\bookref{3.4.39}), 而由对角线的可表性 (定理~\bookref{4.2.1}(2)), 后者是代数空间. 特别地, $\mathcal X$ 具有仿射对角线当且仅当对于每个概形 $S$ 和 $a,b\in\mathcal X(S)$, $\underline{\Isom}_{\mathcal X(S)}(a,b)$ 都仿射. 每个具有仿射或拟仿射对角线的代数叠必然拟分离.
\begin{remark}\label{src:4.3.12}
对于概形的态射, 上述分离性的定义与通常的概念一致, 因为概形态射的对角线是闭浸入当且仅当它固有.
\end{remark}
\begin{remark}\label{src:4.3.13}
拟分离的 Deligne--Mumford 叠具有有限且约化的稳定子群 (见习题~\bookref{4.2.8}).
\end{remark}
\begin{lemma}\label{src:4.3.14}
设 $S$ 是仿射概形, $G\to S$ 是作用于 $S$ 上代数空间 $U$ 的光滑仿射群概形. 若 $U$ 具有仿射对角线 (相应地, 拟仿射对角线), 则 $[U/G]$ 亦然.
\end{lemma}
\begin{proof}
回忆, 定理~\bookref{4.1.10}~已经证明 $[U/G]$ 是代数叠. 对角线的可表性 (定理~\bookref{4.2.1}(2)) 蕴涵 $[U/G]\to[U/G]\times_S[U/G]$ 是可表态射. 对习题~\bookref{3.4.37}~中的 Cartesian 图
\[
\begin{tikzcd}
G\times_S U\arrow[r]\arrow[d]\arrow[dr,phantom,"\square"]&U\times_S U\arrow[d]\\
{[U/G]}\arrow[r]&{[U/G]\times_S[U/G]}
\end{tikzcd}
\]
应用光滑下降 (命题~\bookref{4.3.6}), 只需证明 $G\times_S U\to U\times_S U$ 仿射 (相应地, 拟仿射). 由于 $G$ 仿射, 复合 $G\times_S U\to U\times_S U\xrightarrow{p_2}U$ 亦然. 另一方面, 由于 $U$ 具有仿射对角线 (相应地, 拟仿射对角线), $p_2:U\times_S U\to U$ 也具有此性质. 态射的消去性质蕴涵 $G\times_S U\to U\times_S U$ 仿射 (相应地, 拟仿射).
\end{proof}
大多数模问题都满足仿射对角线条件 (但例如 $\mathcal M_1$ 除外).
\begin{example}\label{src:4.3.15}
模叠 $\mathcal M_g$ 和 $\cBun_{r,d}(C)$ 具有仿射对角线, 因而拟分离. 对于 $\mathcal M_g$, 结论由该引理及定理~\bookref{4.1.17}~中的商呈示 $\mathcal M_g=[H'/\PGL_{5g-5}]$ 得到, 因为 $H'$ 是射影 Hilbert 概形的局部闭子概形. 稍后将证明 $\mathcal M_g$ 分离 (命题~\bookref{6.5.20}).

类似地, 在定理~\bookref{4.1.23}~中, 我们将 $\cBun_{r,d}(C)$ 的每个拟紧开子叠表为商叠 $[Q'/\GL_N]$, 其中 $Q'$ 是射影 Quot 概形的开子概形. 为证明 $\cBun_{r,d}(C)$ 具有仿射对角线, 只需证明对角线沿态射 $\Spec A\to\cBun_{r,d}(C)\times\cBun_{r,d}(C)$ 作基变换后仿射. 但这样的态射通过某个拟紧开子叠 $\mathcal U\subseteq\cBun_{r,d}(C)$ 的 $\mathcal U\times\mathcal U$ 分解, 而我们知道 $\mathcal U$ 具有仿射对角线.
\end{example}
\begin{example}\label{src:4.3.16}
非分离并 $\mathbb A^\infty\cup_{\mathbb A^\infty\setminus0}\mathbb A^\infty$ 是非拟分离概形的典型例子. 对于代数空间和代数叠, 非拟紧群概形的作用还会带来其他病态现象. 例如, $[\mathbb A^1/\mathbb Z]$ 是非拟分离的代数空间 (见例~\bookref{4.9.36}), 而 $B\mathbb Z$ 是非拟分离的代数叠 (见例~\bookref{4.9.36}).
\end{example}
\subsection{叠的拓扑空间}\label{src:4.3.4-subsection}
我们可以为每个代数叠 $\mathcal X$ 赋予拓扑空间 $|\mathcal X|$.
\begin{definition}[代数叠的拓扑空间]\label{src:4.3.17}
若 $\mathcal X$ 是代数叠, 则定义 $\mathcal X$ 的\newterm{拓扑空间}{topological space} 为由域值态射 $x:\Spec K\to\mathcal X$ 构成的集合 $|\mathcal X|$: 若存在域扩张 $K_1\to K_3$ 和 $K_2\to K_3$, 使 $x_1|_{\Spec K_3}$ 与 $x_2|_{\Spec K_3}$ 在 $\mathcal X(K_3)$ 中同构, 则态射 $x_1:\Spec K_1\to\mathcal X$ 和 $x_2:\Spec K_2\to\mathcal X$ 在 $|\mathcal X|$ 中视为相同. 子集 $U\subseteq|\mathcal X|$ 是开的, 是指存在开子叠 $\mathcal U\subseteq\mathcal X$, 使 $U=|\mathcal U|\subseteq|\mathcal X|$.
\end{definition}
代数叠的态射 $\mathcal X\to\mathcal Y$ 诱导连续映射 $|\mathcal X|\to|\mathcal Y|$.
\begin{exercise}\label{src:4.3.18}
证明若 $\mathcal X$ 是代数叠, $U\subseteq|\mathcal X|$ 是开子集, 则存在约化闭子叠 $\mathcal Z\hookrightarrow\mathcal X$, 使 $|\mathcal Z|=|\mathcal X|\setminus U$.
\end{exercise}
\begin{example}\label{src:4.3.19}
域 $k$ 上的商叠 $[\mathbb A^1/\mathbb G_m]$ 的拓扑空间由两个点组成, 其中 $\mathbb G_m$ 通过标准数乘作用作用于 $\mathbb A^1$. 两个点的代表为 $x_0:\Spec k\xrightarrow{0}\mathbb A^1\to[\mathbb A^1/\mathbb G_m]$ 和 $x_1:\Spec k\xrightarrow{1}\mathbb A^1\to[\mathbb A^1/\mathbb G_m]$. 特别地, 一般点的嵌入 $\Spec k(x)\to\mathbb A^1\to[\mathbb A^1/\mathbb G_m]$ 等价于 $x_1$.
\end{example}
给定 $x\in|\mathcal X|$, 稳定子群 $G_{x_1}$ 依赖于 $x$ 的代表 $x_1:\Spec k\to\mathcal X$ 的选择, 但其维数与这一选择无关, 我们将此维数记为 $\dim G_x$. 类似地, 光滑, 非分歧, 仿射, 有限和约化等性质也与这一选择无关.
\begin{exercise}[简单]\label{src:4.3.20}
设 $x\in|\mathcal X|$ 是代数叠的点, 有两个代表 $x_1:\Spec k_1\to\mathcal X$ 和 $x_2:\Spec k_2\to\mathcal X$.
\begin{enumerate}[(1)]
\item 证明稳定子群 $G_{x_1}$ 光滑 (相应地, 平展, 非分歧, 仿射, 有限) 当且仅当 $G_{x_2}$ 如此.
\item 证明 $\dim G_{x_1}=\dim G_{x_2}$.
\item 若 $\mathcal X$ 是 Deligne--Mumford 叠, 且 $k_1$ 和 $k_2$ 都代数闭, 证明对应于 $G_{x_1}$ 和 $G_{x_2}$ 的抽象离散群同构.
\end{enumerate}
\end{exercise}
由上述习题, 称 $x\in|\mathcal X|$ 具有光滑 (相应地, 平展, 非分歧, 仿射, 有限) 稳定子是有意义的. 对于 Deligne--Mumford 叠 $\mathcal X$, 定义 $x$ 的\newterm{几何稳定子}{geometric stabilizer} 为离散群 $G=G_{\bar x}(k)$, 其中 $\bar x:\Spec k\to\mathcal X$ 是代表 $x$ 的几何点.

现在可以定义代数叠及其态射的拓扑性质.
\begin{definition}\label{src:4.3.21}
若 $|\mathcal X|$ 拟紧, 连通或不可约, 则称代数叠 $\mathcal X$ 拟紧, 连通或不可约. 若 $\mathcal X$ 局部 Noether (定义~\bookref{4.3.7}), 拟分离 (定义~\bookref{4.3.11}(2)) 且拟紧, 则称它是 Noether 的.
\end{definition}
\begin{exercise}[简单]\label{src:4.3.22}
证明代数叠 $\mathcal X$ 拟紧当且仅当存在来自仿射概形的光滑呈示 $\Spec A\to\mathcal X$; 并证明拟分离的代数叠 $\mathcal X$ 是 Noether 的当且仅当存在光滑呈示 $\Spec A\to\mathcal X$, 其中 $A$ 是 Noether 环.
\end{exercise}
\begin{example}\label{src:4.3.23}
模叠 $\mathcal M_g$ 是 Noether 的, 特别地, 它拟紧. 利用定理~\bookref{4.1.17}~中的商呈示 $\mathcal M_g=[H'/\PGL_{5g-5}]$, 这由上述习题得到. 虽然 $\cBun_{r,d}(C)$ 拟分离且局部 Noether, 但它不拟紧, 因而不是 Noether 的. 见 \S\bookref{9.2}.
\end{example}
\begin{exercise}[简单]\label{src:4.3.24}
\begin{enumerate}[(a)]
\item 证明代数叠态射 $\mathcal X\to\mathcal Y$ 满射当且仅当 $|\mathcal X|\to|\mathcal Y|$ 满射.
\item 证明若 $\mathcal X\to\mathcal Y$ 和 $\mathcal Y'\to\mathcal Y$ 是代数叠的态射, 则 $|\mathcal X\times_{\mathcal Y}\mathcal Y'|\to|\mathcal X|\times_{|\mathcal Y|}|\mathcal Y'|$ 满射.
\end{enumerate}
\end{exercise}
\begin{exercise}[简单]\label{src:4.3.25}
若 $\mathcal X$ 是拟紧且局部 Noether 的代数叠, 证明 $|\mathcal X|$ 是 Noether 拓扑空间.
\end{exercise}
\begin{exercise}[中等]\label{src:4.3.26}
若 $U\to\mathcal X$ 是代数叠的光滑呈示, 证明 $|U|\to|\mathcal X|$ 浸没, 即它满射且 $|\mathcal X|$ 带有商拓扑.

\emph{提示:} 若 $V\subseteq|\mathcal X|$ 是子集, 它在 $U$ 中的原像 $U'$ 是开的, 则将 $\mathcal V\subseteq\mathcal X$ 定义为满足 $U'_T=U_T$ 的对象 $T\to\mathcal X$ 构成的子叠.
\end{exercise}
\begin{exercise}[简单]\label{src:4.3.27}
代数叠的泛开态射可以利用定义~\bookref{4.3.2}(1) 定义, 因为它们在源和靶上光滑局部. 例如, 局部有限表现的忠实平坦态射是泛开的. 若 $f:\mathcal X\to\mathcal Y$ 是代数叠的泛开态射, 证明 $f(|\mathcal X|)\subseteq|\mathcal Y|$ 是开的. 推出对于每个代数叠态射 $\mathcal Y'\to\mathcal Y$, 映射 $|\mathcal X\times_{\mathcal Y}\mathcal Y'|\to|\mathcal Y'|$ 是开的.
\end{exercise}
\begin{definition}\label{src:4.3.28}
代数叠态射 $\mathcal X\to\mathcal Y$ 拟紧, 是指对于每个态射 $\Spec B\to\mathcal Y$, 纤维积 $\mathcal X\times_{\mathcal Y}\Spec B$ 拟紧. 称 $\mathcal X\to\mathcal Y$ 有限型, 如果它局部有限型且拟紧.
\end{definition}
\begin{example}\label{src:4.3.29}
模叠 $\mathcal M_g$ 在 $\mathbb Z$ 上有限型. 另一方面, $\cBun_{r,d}(C)$ 在 $k$ 上局部有限型, 但不是有限型的.
\end{example}
\begin{remark}\label{src:4.3.30}
拟紧态射 $\mathcal X\to\mathcal Y$ 在拓扑空间上诱导拟紧映射 $|\mathcal X|\to|\mathcal Y|$. 若 $\mathcal Y$ 拟分离, 则逆命题成立, 但一般情形并非如此, 例如 $\Spec k\to B\mathbb Z$ (见例~\bookref{4.9.36}).
\end{remark}
\begin{exercise}[有挑战性]\label{src:4.3.31}
\begin{enumerate}[(a)]
\item 设 $f:\mathcal X\to\mathcal Y$ 是代数叠的拟紧态射. 对于 $y\in|\mathcal Y|$, 证明 $y\in f(|\mathcal X|)$ 当且仅当存在 $x\in|\mathcal X|$ 及特殊化 $f(x)\rightsquigarrow y$.
\item 证明代数叠的开态射 $f:\mathcal X\to\mathcal Y$ (即 $|\mathcal X|\to|\mathcal Y|$ 是开的) 满足下述提升性质: 若 $x\in|\mathcal X|$ 是点, 则每个特殊化 $y'\rightsquigarrow f(x)$ 都提升为特殊化 $x'\rightsquigarrow x$. 证明对于有限表现态射, 逆命题成立.
\item 若 $\mathcal X$ 是拟分离的代数叠, 证明 $|\mathcal X|$ 是\newterm{清醒拓扑空间}{sober topological space}, 即每个不可约闭子集都有唯一的一般点.
\item 将 Chevalley 判据推广到代数叠: 若 $f:\mathcal X\to\mathcal Y$ 是代数叠的有限表现态射, 则像 $f(|\mathcal X|)\subseteq|\mathcal Y|$ 是可构造的.
\end{enumerate}
另见 \cite[\S$5.7$--$5.9$]{LMB00} 和 \cite[Tags 0DQQ 和 0GVY]{SP}.
\end{exercise}
\begin{exercise}[一般平坦性]\label{src:4.3.32}
将一般平坦性 (\bookref{A.2.13}) 从概形态射推广到代数叠: 若 $\mathcal X\to\mathcal Y$ 是代数叠的有限型态射, 且 $\mathcal Y$ 约化, 则存在稠密开子叠 $\mathcal U\subseteq\mathcal Y$, 使基变换 $\mathcal X_{\mathcal U}\to\mathcal U$ 平坦且有限表现.
\end{exercise}
\begin{exercise}[技术性]\label{src:4.3.33}
将命题~\bookref{A.1.1}~给出的局部有限表现态射的刻画推广到代数叠: 代数叠的态射 $f:\mathcal X\to\mathcal Y$ 局部有限表现, 当且仅当对于 $\mathcal Y$ 上仿射概形的每个有向系统 $\{\Spec A_\lambda\}_{\lambda\in I}$, 自然映射
\[
\operatorname{colim}_\lambda\Mor_{\mathcal Y}(\Spec A_\lambda,\mathcal X)
\longrightarrow\Mor_{\mathcal Y}(\Spec(\operatorname{colim}_\lambda A_\lambda),\mathcal X)
\]
是范畴等价.
\end{exercise}
\subsection{拟有限态射}\label{src:4.3.5-subsection}
概形的态射局部拟有限, 是指它局部有限型且每个纤维离散, 即带有离散拓扑. 由于此性质在源和靶上平展局部, 可以利用定义~\bookref{4.3.2}~将它推广到代数空间的态射.
\begin{definition}\label{src:4.3.34}
\begin{enumerate}[(1)]
\item 代数叠的可表态射 $\mathcal X\to\mathcal Y$ 局部拟有限, 是指对于每个来自概形的态射 $T\to\mathcal Y$, 代数空间 $\mathcal X\times_{\mathcal Y}T$ 在 $T$ 上局部拟有限.
\item 代数叠的态射 $\mathcal X\to\mathcal Y$ 局部拟有限, 是指它局部有限型, 对角线 $\mathcal X\to\mathcal X\times_{\mathcal Y}\mathcal X$ 局部拟有限, 且对于每个来自域的态射 $\Spec k\to\mathcal Y$, 拓扑空间 $|\mathcal X\times_{\mathcal Y}\Spec k|$ 离散.
\item 代数叠的态射 $\mathcal X\to\mathcal Y$ 拟有限, 是指它局部拟有限且拟紧.
\end{enumerate}
\end{definition}
为理解定义~\bookref{4.3.34}(2) 中的条件, 回忆对角线 $\mathcal X\to\mathcal X\times_{\mathcal Y}\mathcal X$ 总是可表态射 (习题~\bookref{4.2.4}). 对角线拟有限 (相应地, 局部拟有限), 当且仅当对于像为 $y\in\mathcal Y(k)$ 的每个域值点 $x\in\mathcal X(k)$, 稳定子群的诱导映射的核 $\ker(G_x\to G_y)$ 有限 (相应地, 离散); 见习题~\bookref{4.2.7}. 特别地, 若 $\mathcal Y$ 是概形, 则对角线拟有限当且仅当 $\mathcal Y$ 的所有稳定子都有限. 例如, 若 $G$ 是域 $k$ 上的有限群概形 (例如 $\mu_p$), 则 $BG\to\Spec k$ 拟有限. 另一方面, 尽管 $|B\mathbb G_m|$ 只有一个点, $B\mathbb G_m\to\Spec k$ 仍不拟有限.
\begin{exercise}[简单]\label{src:4.3.35}
证明代数叠的有限型态射 $f:\mathcal X\to\mathcal Y$ 拟有限, 当且仅当 $|\mathcal X|\to|\mathcal Y|$ 具有有限纤维, 且对于每个域值点 $x\in\mathcal X(k)$, 映射 $\Aut_{\mathcal X(k)}(x)\to\Aut_{\mathcal Y(k)}(f(x))$ 具有有限余核.
\end{exercise}
稍后将证明每个可表, 拟有限且分离的态射都拟仿射 (命题~\bookref{5.5.5}).
\subsection{平展与非分歧态射}\label{src:4.3.6-subsection}
我们要求平展和非分歧态射是相对 Deligne--Mumford 的.
\begin{definition}\label{src:4.3.36}
$\Sch_{\text{ét}}$ 上的叠态射 $\mathcal X\to\mathcal Y$ 是\newterm{相对 Deligne--Mumford}{relatively Deligne--Mumford} 的, 是指对于每个来自概形的态射 $T\to\mathcal Y$, 纤维积 $\mathcal X\times_{\mathcal Y}T$ 是 Deligne--Mumford 叠.
\end{definition}
在推论~\bookref{4.6.5}~中将看到, 相对 Deligne--Mumford 态射由其对角线的非分歧性刻画. 满足较弱条件, 即对角线 $\mathcal X\to\mathcal X\times_{\mathcal Y}\mathcal X$ 局部拟有限的态射 $\mathcal X\to\mathcal Y$, 称为\newterm{拟 DM}{quasi-DM} 态射 \cite[Tag 04YW]{SP}.

对于概形的态射, 平展性和非分歧性在源上平展局部, 在靶上光滑局部. 因而这些概念推广到相对 Deligne--Mumford 态射.
\begin{definition}\label{src:4.3.37}
代数叠的态射 $\mathcal X\to\mathcal Y$ 平展 (相应地, 非分歧), 是指它是相对 Deligne--Mumford 的, 且对于每个光滑呈示 $V\to\mathcal Y$ 和平展呈示 $U\to\mathcal X\times_{\mathcal Y}V$, 诱导的概形态射 $U\to V$ 平展 (相应地, 非分歧).
\end{definition}
态射平展当且仅当它光滑且非分歧, 态射非分歧当且仅当其对角线平展. 这些结论由概形态射的相应事实 (定理~\bookref{A.3.2}~和~\bookref{A.3.3}) 得到; 注意对于概形的态射, 对角线平展当且仅当它是开浸入.

虽然平展态射光滑且局部拟有限, 但逆命题并不成立. 例如, 在特征为 $p$ 的域 $k$ 上, 映射 $B\mu_p\to\Spec k$ 光滑且拟有限, 但不平展, 因为 $B\mu_p$ 不是 Deligne--Mumford 叠 (见习题~\bookref{7.3.12}). 类似地, 平展态射光滑且相对维数为 $0$, 但逆命题仍然不成立, 例如特征为 $p$ 时的 $B\mu_p\to\Spec k$, 或任意特征时的 $[\mathbb A^1/\mathbb G_m]\to\Spec k$.
\begin{exercise}[技术性]\label{src:4.3.38}
对于代数叠态射 $f:\mathcal X\to\mathcal Y$, 证明下列条件等价:
\begin{enumerate}[(1)]
\item 对角线 $\Delta_f:\mathcal X\to\mathcal X\times_{\mathcal Y}\mathcal X$ 非分歧 (相应地, 分离, 拟分离),
\item 相对惯性叠 $I_{\mathcal X/\mathcal Y}\to\mathcal X$ 非分歧 (相应地, 分离, 拟分离), 且
\item 双重对角线 (等价地, 相对惯性叠的单位截面)
\[
\mathcal X\longrightarrow I_{\mathcal X/\mathcal Y}
=\mathcal X\times_{\mathcal X\times_{\mathcal Y}\mathcal X}\mathcal X
\]
是开浸入 (相应地, 闭浸入, 拟紧).
\end{enumerate}
\emph{提示:} 将命题~\bookref{B.1.9}(4) 对非分歧 (相应地, 分离, 拟分离) 群概形的刻画推广到群代数空间.
\end{exercise}
